train 拟合模型，add 将行编码并按列表排序，search 使用完成的索引。缓存命中跳过拟合，但不跳过编码，因此冷构建与缓存搜索的成本不同。

\begin{figure}[!htbp]
\centering
% Audited lifecycle; no proposed storage or distributed subsystem.
\begin{tikzpicture}[
 font=\small\linespread{1}\selectfont,
 box/.style={draw=ThesisNavy!65,rounded corners=2pt,align=center,
   text width=3.65cm,minimum height=1.9cm,inner sep=5pt},
 gpu/.style={box,fill=ThesisNavy!7},
 arr/.style={-{Latex[length=2mm]},draw=ThesisNavy!80,thick}]
\node[box] (host) at (0,0) {Host input vectors\\Training samples\\Database batches};
\node[box] (model) at (4.8,0) {Host model\\QR rotation $\Pi$\\Primary codebook};
\node[box] (cache) at (9.6,0) {Host model cache\\$\Pi$ and codebooks\\IVF centroids};
\node[gpu] (train) at (0,-3.6) {GPU training\\Rotate samples\\Fit IVF centroids\\Fit residual model};
\node[gpu] (add) at (4.8,-3.6) {GPU bulk build\\1 or 2 passes\\Assign / encode\\Reorder / pack\\Floats $\leq768$ MiB};
\node[gpu] (index) at (9.6,-3.6) {GPU-resident index\\Offsets and IDs\\Primary / residual codes\\Corrections};
\node[box] (query) at (0,-7.2) {Host query batch\\Pinned staging\\Transfer to GPU};
\node[gpu] (search) at (4.8,-7.2) {GPU query pipeline\\Rotate; probe; scan\\Refine; final top-$k$\\Workspace / graph};
\node[box] (result) at (9.6,-7.2) {Host query results\\IDs and scores\\Transfer from GPU};
\draw[arr] (host) -- node[left,font=\footnotesize] {samples} (train);
\draw[arr] (model.south west) -- (train.north east);
\draw[arr] (train) -- node[above,font=\footnotesize] {model} (add);
\draw[arr] (host.west) -- ++(-.5,0) |- (4.8,-5.15)
 node[pos=.8,below,font=\footnotesize] {database batches} -- (add.south);
\draw[arr] (model) -- (add);
\draw[arr] (add) -- (index);
\draw[<->,dashed,draw=ThesisNavy!80,thick]
 (train.north east) -- node[pos=.8,above,sloped,font=\footnotesize,fill=white,inner sep=1pt]
 {save / load model} (cache.south west);
\draw[arr] (query) -- (search);
\draw[arr] (search) -- (result);
\draw[arr] (index.south) -- ++(0,-.6) -| (search.north);
\end{tikzpicture}

\caption[构建与查询架构]{实现中的构建与查询边界。主机工作包括抽样、正交矩阵构造和解析主码本生成；评测配置在 GPU 上拟合 IVF 与残差模型。缓存保存训练状态，完整数据库表示则位于设备上。虚线表示模型复用。暂存上限不包含其他临时空间。}
\label{fig:4-architecture}
\end{figure}

主机准备样本、旋转和主码本；GPU 执行旋转，拟合 IVF／残差模型，并将有界 batch 编码为列表顺序。

模型保存旋转、码本和质心；数据库状态保存 offsets、ID、编码和校正项。CUDA Graph 重放已记录的 GPU 操作，并引用复用的查询缓冲。

一次成功的 \texttt{add} 固定物理位置。后续 add 被拒绝，更新需要重建。可变／事务性发布尚未实现。

主机 ID／分数及调用计时遵循 \cref{eq:5-refined-constant,sec:3-3}。
